\chapter{Arquitectura Computacional y Trazabilidad}
\label{chap:arquitectura}

\section{Motivación de la arquitectura}
Una tesis experimental de nivel de magíster no puede depender de resultados aislados obtenidos manualmente. La validez de las conclusiones exige que cada número reportado pueda relacionarse con un script, una configuración, una semilla aleatoria y un archivo de salida. Por esta razón, el trabajo se organizó como un \textit{toolkit} modular de experimentación, donde la reconstrucción de canales EEG se ejecuta mediante componentes desacoplados: carga de datos, construcción de grafos, simulación de canales faltantes, interpolación, evaluación y consolidación estadística.

La arquitectura persigue tres objetivos. Primero, permitir incorporar nuevos métodos de interpolación sin reescribir el pipeline completo. Segundo, facilitar corridas exhaustivas y repetibles con variaciones controladas de dataset, escenario de pérdida e hiperparámetros. Tercero, producir artefactos legibles tanto por scripts de análisis como por el documento de tesis: CSV, tablas en \LaTeX, figuras y reportes de trazabilidad.

\section{Estructura lógica del pipeline}
La Figura~\ref{fig:pipeline_tikz} resume el flujo experimental. Cada bloque genera salidas persistentes y versionables. Esta decisión es importante porque las corridas extensas pueden fallar por límites de memoria, tiempo o disponibilidad de datos; al guardar resultados intermedios, una ejecución puede reanudarse sin repetir trabajo ya completado.

\begin{figure}[H]
\centering
\begin{tikzpicture}[
  node distance=1.55cm,
  box/.style={rectangle, draw=hermesblue, rounded corners, minimum height=0.9cm, minimum width=2.35cm, align=center, font=\small},
  arrow/.style={-{Stealth[length=2.4mm]}, thick, draw=hermesgray}
]
\node[box, fill=blue!8] (data) {Carga EEG\\datasets};
\node[box, fill=green!8, right=of data] (graph) {Construcción\\de grafo};
\node[box, fill=orange!10, right=of graph] (mask) {Máscaras\\faltantes};
\node[box, fill=purple!8, below=of graph] (interp) {Interpolación\\GSP/TV/baseline};
\node[box, fill=red!8, right=of interp] (metrics) {Métricas\\MAE--LSD};
\node[box, fill=gray!12, left=of interp] (optuna) {Optuna\\búsqueda};
\node[box, fill=yellow!18, below=of interp] (report) {Tablas, figuras\\y reportes};
\draw[arrow] (data) -- (graph);
\draw[arrow] (graph) -- (mask);
\draw[arrow] (mask) |- (interp);
\draw[arrow] (optuna) -- (interp);
\draw[arrow] (interp) -- (metrics);
\draw[arrow] (metrics) |- (report);
\draw[arrow] (report) -| (optuna);
\end{tikzpicture}
\caption{Arquitectura lógica del pipeline experimental. La optimización de hiperparámetros y la generación de reportes se integran como un ciclo reproducible, no como una etapa manual posterior.}
\label{fig:pipeline_tikz}
\end{figure}

\section{Convención de artefactos}
Los resultados se organizaron en familias de artefactos. Los archivos \code{*_raw.csv} conservan observaciones por ensayo o configuración; los archivos \code{*_summary.csv} agregan medidas descriptivas; los reportes \code{*.md} documentan decisiones, advertencias y estado de cierre. Las tablas \LaTeX{} exportadas se usan en el manuscrito, pero no reemplazan a los CSV de origen. Esta separación evita que una tabla formateada se convierta en la única fuente de verdad.

\begin{table}[H]
\centering
\caption{Tipos de artefactos generados por el pipeline y función dentro de la tesis.}
\label{tab:artifact_types}
\begin{tabular}{p{0.22\textwidth}p{0.28\textwidth}p{0.36\textwidth}}
\toprule
Tipo & Ejemplo & Función \\
\midrule
Datos crudos & \code{*_raw.csv} & Preservan observaciones por ensayo, semilla, dataset y método. \\
Resumen & \code{*_summary.csv} & Agregan medias, medianas, desviaciones estándar e intervalos. \\
Ranking & \code{*_ranking*.csv} & Ordenan métodos o familias bajo una métrica objetivo. \\
Figura & \code{*.png}, \code{*.pdf} & Permiten inspección visual de degradación, ERP, PSD y trade-offs. \\
Reporte & \code{*.md} & Documentan decisiones editoriales, reproducibilidad, advertencias y limitaciones. \\
Tabla tesis & \code{*.tex} & Versiones tipográficas derivadas de los CSV, no fuentes primarias. \\
\bottomrule
\end{tabular}
\end{table}

\section{Política de normalización}
Una decisión metodológica importante fue separar corridas normalizadas y no normalizadas. Mezclar métricas de MAE o RMSE obtenidas bajo escalas distintas produce comparaciones inválidas, especialmente al cruzar datasets con amplitudes y densidades espaciales diferentes. Por ello, cada corrida debe registrar el campo \code{normalization}; si el campo difiere entre dos artefactos, estos no deben agregarse en una misma inferencia. Esta política se refleja en la documentación del repositorio y se mantiene como restricción de lectura para los resultados de esta tesis.

\section{Trazabilidad editorial}
La redacción del documento se apoya en una regla simple: todo resultado declarado debe vincularse a un artefacto. La Tabla~\ref{tab:claim_traceability} presenta ejemplos de esta trazabilidad. En la versión final de una tesis institucional, esta tabla puede moverse al apéndice; aquí se mantiene en el cuerpo porque define cómo deben interpretarse las afirmaciones cuantitativas.

\begin{table}[H]
\centering
\caption{Ejemplos de trazabilidad entre afirmaciones, artefactos y secciones.}
\label{tab:claim_traceability}
\begin{tabular}{p{0.30\textwidth}p{0.30\textwidth}p{0.25\textwidth}}
\toprule
Afirmación & Artefacto fuente & Sección asociada \\
\midrule
Comparación final por familia & \code{b3\_rep01\_final\_table\_overall.csv} & Resultados globales \\
Significancia TRSS vs BGSRP & \code{b3\_stat02\_significance.csv} & Validación estadística \\
Estado parcial BGSRP & \code{ins13\_strict\_status.md} & Discusión y limitaciones \\
Recursos computacionales & \code{b3\_b4\_compute\_resources.md} & Reproducibilidad \\
Figuras ERP/PSD & \code{results\_optuna\_final/} y \code{results\_optuna\_baselines/} & Fidelidad fisiológica \\
\bottomrule
\end{tabular}
\end{table}

\section{Reproducibilidad práctica}
La reproducibilidad práctica no se limita a publicar código. También requiere especificar ambiente de ejecución, semillas, parámetros de entrada y criterios de exclusión. En esta tesis se utiliza \code{random\_state=42} cuando corresponde, se documentan configuraciones de Optuna, se conservan advertencias clasificadas y se registran estados \textit{GO/NO-GO}. Esta última distinción evita presentar iteraciones exploratorias o fallidas como evidencia confirmatoria.
